The executable GPU evidence was obtained on 13 November 2026 on an
NVIDIA GeForce RTX 4090.  The recorded environment reports
25,250,627,584 bytes of device memory and 23,960,354,816 initially free
bytes.  It uses Python 3.12.3, cuQuantum Python 26.11.0, cuStateVec
1.15.1, cuStabilizer 0.5.0, and CuPy 14.2.0.  The original record is
\path{extension/quantum/results/gpu-results.json}; it contains nine
successful experiments, the numerical outputs, timestamps, package
versions, and the SHA256 of the executed script.  The source and the
record are retained separately so that a later, larger test does not
silently replace earlier evidence.

\paragraph{Dense state-vector results.}
A two-qubit Bell-state calculation first checked amplitude ordering and
$XX,ZZ$ expectations.  The surface-code runs then executed the emitted
Hadamard/CNOT encoder using actual cuStateVec matrix applications.  They
computed all star and face expectations, injected an $X$ error on a
specified edge, and compared every resulting expectation with the
predicted incident-face signs.  One thirty-qubit run appended a fresh
ancilla and four data-to-ancilla CNOTs to extract one face syndrome.
This final experiment explicitly stored a thirty-one-qubit state.

\begin{center}
\begin{tabular}{lrrrr}
\hline
Data and ancillas & Raw state & Encoder gates & Preparation (s) & Timed total (s)\\
\hline
$18+0$ & 2 MiB & 42 & 0.000312 & 0.000537\\
$30+0$ & 8 GiB & 74 & 0.918804 & 1.204211\\
$30+1$ & 16 GiB & 74 & 1.838463 & 2.528538\\
\hline
\end{tabular}
\end{center}

All three cases use complex64 amplitudes.  Their recorded norm squared
is approximately $0.9999997616$ and their maximum pre-error stabilizer
expectation discrepancy is $2.385\times10^{-7}$.  The extra ancilla has
$Z$ expectation approximately $-0.9999997616$, as predicted by the
injected error.  The routine evaluates these expectations directly; it
does not simulate a finite collection of measurement shots.  The listed
raw-state sizes are exact allocations implied by the precision and
qubit count.  The required gate workspace was zero in these runs.

Preparation time covers encoder application and synchronization.  The
timed total begins after construction of the state-vector object and
includes preparation, the reported expectation checks, and the inserted
error and extraction gates.  It excludes state allocation and final
cleanup.  These are measured intervals for this particular implementation,
not complete application launch times or portable hardware benchmarks.

\paragraph{Training a circuit that has an exact analytic baseline.}
The eighteen-qubit neural experiment executed the encoded preparation,
rotations, fan-in CNOTs, and output expectation on the GPU.  Each gradient
component was obtained from the parameter-shift expectation evaluations;
the analytic product formula was used as a validation reference, not as
a replacement for the quantum gates.  Twenty-four updates at learning
rate $0.2$ changed the initial mean squared error
$0.21987940665826008$ to $0.0027618471230164243$, a reduction by a
factor of approximately $79.6$.  The recorded final angles are
\[
 (0.1918860335,-0.1237818552,0.1522945504,-0.1791698376).
\]
The largest discrepancy from the exact output formula along the recorded
training history was $3.093\times10^{-7}$.  All sixteen final truth-table
entries had the correct sign, including the twelve inputs omitted from
the updates.  The training and evaluation interval was approximately
$0.045774$ seconds, excluding the initial encoded-state construction.

The initial angles already give $C(\theta)>0$.  Consequently the
thresholded sign classifier is correct on all sixteen inputs before
training as well as afterwards.  The measured improvement is the
calibration of expectation magnitude and mean squared error; it is not
an increase from poor classification accuracy to perfect classification.
The magnitude changes from approximately $0.531087$ to $0.947447$.
If a single ideal measurement were used to predict parity, the exact
Born-rule success probability $(1+C)/2$ would accordingly change from
approximately $0.76554$ to $0.97372$.  These probabilities are deductions
from computed expectations, not empirically sampled success frequencies.

\paragraph{Pauli frames and the reference-state assumption.}
For the larger experiments, cuStabilizer stores batched Pauli errors
relative to a known ideal encoded reference.  It does not store the
$2^n$ amplitudes of that reference.  The mathematical encoder theorem
supplies an ideal state with every specified check equal to $+1$.
The frame experiments assume that state and propagate its errors through
ideal syndrome extraction.  They do not execute the large encoder or
construct a generic stabilizer tableau for an arbitrary unknown state.

Each experiment uses 32,768 error realizations with seed 20261113.
Every data qubit receives an independent $X$ error with probability
$0.001$ and an independent $Z$ error with the same probability.  Both
can occur, giving the corresponding $Y$ error up to an irrelevant phase.
Syndrome ancillas and extraction gates are noiseless.  A star ancilla
is reset, Hadamard transformed, used as the control for its incident
CNOTs, Hadamard transformed again, and measured.  A face ancilla receives
CNOTs controlled by the boundary data qubits and is then measured.
Measurement randomization is disabled because the stored output is the
error-induced sign change from the assumed deterministic ideal outcome.

Let $x,z\in\mathbb F_2^E$ be the data error masks.  The exact predicted
syndrome is
\[
 s_X=H_Xz,\qquad s_Z=H_Zx.
\]
With these noiseless, freshly reset ancillas, the extraction circuit
preserves the data error masks.  Hence the implementation can compare
its measured sign-change bits with sparse binary products computed from
the returned data frames.  It does so for every stored syndrome byte,
not merely for a random sample.  Every comparison in the recorded runs
has zero mismatches.  This checks the consistency of the extraction
circuit, qubit indexing, packed-bit representation and binary incidence
calculation.  It does not test a recovery algorithm or estimate a logical
failure probability.

\begin{center}
\begin{tabular}{lrrrr}
\hline
Case & Data & Ancillas & Apply interval (s) & Timed total (s)\\
\hline
Genus two & 32 & 30 & 0.004756 & 3.025362\\
Torus $L=32$ & 2,048 & 2,048 & 0.005072 & 0.013088\\
Torus $L=128$ & 32,768 & 32,768 & 0.073878 & 0.165669\\
Torus $L=512$ & 524,288 & 524,288 & 1.147181 & 2.693085\\
\hline
\end{tabular}
\end{center}

Here the apply interval is a synchronized wall-clock interval around
noise and extraction circuit application, including host dispatch.  It
is not a CUDA-event measurement of kernel time alone.  The timed total
includes simulator setup and exhaustive syndrome verification, but
excludes topology generation before entry into the timer and final
cleanup afterwards.  First-use compilation and initialization make the
genus-two total much larger than its apply interval.  No cross-backend
speedup is inferred from these differently scoped measurements.

For the $L=512$ experiment, all 2,147,483,648 packed syndrome bytes were
checked.  There were 68,489,748 violated checks over all shots, giving
an aggregate syndrome rate of approximately $0.00398663$.  For a
weight-four check under independent relevant Pauli errors of probability
$p$, the exact odd-parity probability is
\[
 \frac{1-(1-2p)^4}{2}.
\]
At $p=0.001$ this is approximately $0.00398802$.  This agreement is a
useful noise-model sanity check.  Correlations among neighboring checks
and their global parity constraints mean that the aggregate count should
not be interpreted as independent Bernoulli samples at every check.

\paragraph{Memory accounting and the scope of a capacity claim.}
For $S$ shots, the stored row stride in bytes is
$4\lceil S/32\rceil$.  With $N$ total data-plus-ancilla qubits and $A$
measured ancillas, the implementation estimates the explicit packed bit
tables as
\[
 B=(2N+A)\,4\lceil S/32\rceil.
\]
The factor two accounts for the $X$ and $Z$ frame components.  This is
an analytic allocation estimate, not an instrumented peak-memory
measurement.  The first run's $L=512$ case has $B=10$ GiB.  Before
allocation, the program requires $1.25B$ to fit within available device
memory after retaining a fixed 2 GiB reserve.  This policy leaves room
for temporary arrays and the desktop; it does not establish a universal
memory bound for every library configuration.

A separate followup script, \path{extension/quantum/gpu_capacity.py},
retains the original experiment source and result unchanged.  It hashes
its own source, the imported experiment source, the topology source,
all archived case JSON files, and the original result.  Its capacity
rule fixes 32,768 shots and chooses the largest multiple of 32 satisfying
\[
 1.25\,(10L^2)(4096)\leq
       \text{free device bytes}-2\,\mathrm{GiB}.
\]
The resulting test is a bounded workload selected under stated memory
assumptions.  A different shot count, error representation, precision,
number of observables, or available-memory state changes the achievable
qubit count.  The dense thirty-one-qubit result and the much larger
frame results therefore answer different computational questions.

\paragraph{Separate capacity followup.}
The followup ran from 10:58:12 to 10:58:17 UTC on the same date and
device.  It successfully exercised the genus-three code with 46 data
qubits and 42 syndrome ancillas, and then selected $L=640$ from the
observed free-memory value of 23,922,606,080 bytes.  The latter case has
819,200 data qubits and 819,200 ancillas, for a total of 1,638,400
qubits in each of 32,768 Pauli frames.  Its estimated packed tables
occupy 16,777,216,000 bytes, or 15.625 GiB.  The budget including the
25 percent overhead allowance and 2 GiB reserve is 23,119,003,648
bytes, below the recorded available amount.

\begin{center}
\begin{tabular}{lrrrr}
\hline
Case & Apply (s) & Timed body (s) & Full call (s) & Checked bytes\\
\hline
Genus three & 0.002198 & 0.081878 & 0.209550 & 172,032\\
Torus $L=640$ & 1.791833 & 4.327947 & 4.813268 & 3,355,443,200\\
\hline
\end{tabular}
\end{center}

Both cases had zero syndrome-byte mismatches.  The larger test observed
107,029,374 violated checks, an aggregate rate of approximately
$0.00398715$.  The additional full-call timer includes topology
generation and cleanup; it still excludes Python startup, library
imports, and output-file serialization.  These boundaries make the
measurements reproducible without calling a host wall-clock interval
a pure kernel time.

The separate evidence file is
\path{extension/quantum/results/capacity-results.json}.  Its source and
case hashes agree with the archived local inputs, and it records that
those inputs were unchanged during execution.  A byte-exact copy of the
original first-run JSON is retained as
\path{extension/quantum/results/gpu-results-remote-original.json}; the
ordinary first-run JSON has identical parsed values but a few different
numeric formatting choices.  This distinction explains their different
file hashes without conflating it with a difference in experiment data.
The largest completed frame workload is therefore the stated $L=640$
case under the stated memory policy.  This observation establishes
neither an absolute maximum qubit count nor the feasibility of arbitrary
non-Clifford computation at that scale.
